\section{Implementation Details}
\label{appendix:imple_detail}

% \paragraph{Implementation Details}
We implement ToG-3 experiments with the following configuration:
\textbf{Data Processing}: Chunk size is set to 1024 tokens with 20-token overlap between consecutive chunks to maintain contextual continuity.
\textbf{Multi-Agent hyperparameter}: Constructor Agent extracts a maximum of 2 knowledge triplets per chunk and employs hierarchical Leiden clustering \citep{traag2019Leiden} with maximum cluster size of 5 for community detection. Retriever Agent retrieves top-5 most relevant nodes using hybrid vector-graph similarity matching. Then, the Reranker reranks the top-2 relevant evidence nodes (or triples) within this retrieved subgraph. Reflector/Responser Agent utilizes the top-2 retrieved passages as context for answer generation.
\textbf{Backend Infrastructure}: LLM service is based on Qwen2.5-32B-Instruct \citep{qwen2.5} deployed with vLLM \citep{vllm} engine using bfloat16 precision and prefix caching enabled and greedy-search generation method, which is more stable than the Qwen3 model in mixed reasoning mode in our task; embeddings are generated using Jina-embeddings-v3 (1024-dimensional) \citep{sturua2024jinaembeddingsv3}; we use jina-reranker-v2~\citep{wang2025jina-reranker-v3} as the reranker model; 
Our server is equipped with 8 A100 40GB cards, AMD EPYC 256-core Processor, 2TB memory, and Ubuntu 20.04.1 system.
and the hybrid vector-graph storage is implemented using Neo4j community edition \footnote{https://neo4j.com/product/community-edition} for efficient knowledge representation and retrieval, see Appendix.\ref{appendix:graph_vis_examples} for visualized graph example.